\section{Empirical implications and policy interpretation}\label{sec:empirical}

The central policy proposition distinguishes the region's \emph{commercialisation incidence} $K_i=M_i+\omega_i\Pi$ from its willingness to pay for restored basic service, $T_i$. Empirically, a counterintuitive screening--incidence inversion would require three separate facts: (i) higher skill complementarity is associated with a higher privately optimal fee and greater exclusion among otherwise eligible firms; (ii) venture launch and competition produce a negative combined regional consumer--incumbent--assigned-founder market effect $K_i$; and (iii) the marginal welfare value of basic service access is sufficiently high that the public-access cost ceiling rises. A descriptive project count cannot verify any of these signs. The primitive inequality \eqref{eq:screen-inversion-band} additionally requires identifying the local appropriation share $\omega_i$, the net incumbent/consumer effect $M_i$, and the ratio $u_i/\Delta_i$; these are economic constructs, not quantities directly recorded in a programme dashboard. The weak-matching limit and uniform-ability assumption are maintained theory restrictions, not empirically established for any jurisdiction.

Although the paper is theoretical, its mechanism can be confronted with data on fees, applicants, matched ventures and downstream markets. The distinction between \emph{conditional model predictions} and \emph{causal estimates} is essential. The following three implications follow from existing equations; they are neither additional general equilibrium theorems nor claimed findings from newly collected data. In particular, posted prices, access to basic services and the availability of successful partners are jointly endogenous outside the stated conditional comparisons.

\paragraph{Implication 1: private fees select participant composition as well as participation volume.}
Hold the other region's moments $(m_j,q_j)$, the service's nonprice characteristics and the local coefficients $U_i>0$ and $T_i>0$ fixed. For a private-only service on its interior cutoff branch, $\ell_i=(f_i-U_i)/T_i$, $m_i=1-\ell_i$ and $q_i=(1-\ell_i^2)/2$. Consequently,
\begin{equation}
 \frac{\partial m_i}{\partial f_i}=-\frac1{T_i},\qquad
 \frac{\partial q_i}{\partial f_i}=-\frac{\ell_i}{T_i},\qquad
 \frac{\partial(q_i/m_i)}{\partial f_i}=\frac1{2T_i}>0,
 \quad 0<\ell_i<1.
 \label{eq:empirical-screen}
\end{equation}
Thus an exogenous fee increase reduces founder participation and the \emph{total} represented skill moment, while raising participants' \emph{average} technical ability. The last effect does not establish an increase in the number of successful projects. A possible empirical test would link the full universe of eligible firms to fee offers, take-up decisions, \emph{pre-service} assessments of technical ability, and subsequent collaborations. After normalising a defensible ex-ante ability measure, one could separately estimate take-up $m_i$, participant mean skill $q_i/m_i$, and the skill moment $q_i$, rather than using registration or match counts as sufficient proxies for innovation quality.

Fee endogeneity is the principal threat: specialists may raise prices because expected demand or client skill improved; targeted discounts may also depend on unobserved ability. A credible design needs plausibly exogenous fee variation (for example, an independently determined tariff change) and a known eligible denominator, not merely a positive correlation between fees and average participant quality. If willingness to pay is poorly aligned with the technical skill relevant to the venture, the model's $z_i=\theta_i$ assumption fails and the positive composition prediction is not warranted. The evidence from European Digital Innovation Hubs that supported client firms can be relatively capable \citep{CarpentierDAddaNepelski2026} reinforces the importance of measuring \emph{who participates}; it does not identify this model's fee-selection mechanism.

\paragraph{Implication 2: public basic access may reduce private fee receipts without eliminating a specialist.}
With an active private provider, a free basic public option and the uniform $z_i=\theta_i$ benchmark, the premium is selected by $\theta_i\geq f_i/\Delta_i$. Its optimal price and take-up share are $f_i^B=\Delta_i/2$ and $D_i^B=1/2$, and gross fee receipts equal $N_i\Delta_i/4$ (equation~\eqref{eq:coexist}). Relative to the private-only fee optimum, with the \emph{other jurisdiction's participation held fixed}, equation~\eqref{eq:revenue-crowding} implies $\mathcal R_i^P>\mathcal R_i^B$. The private provider nevertheless remains viable under co-provision when $F_{P_i}<N_i\Delta_i/4$. These are predictions of a particular service differentiation and cost structure, not a universal claim that public programmes displace private firms.

An evaluation should track separately the universe served by free basic services, paid premium upgrades, actual paid fees (net of vouchers), provider receipts and fixed operating costs, and exits or continued operation. A decline in premium revenue does not imply exit. Comparing different jurisdictions before and after public opening without accounting for changes in specialist quality, local demand, rival services, subsidies and fee policies would not isolate the conditional displacement effect. The prediction of a one-half premium upgrade rate is especially restrictive: it is tied to uniform willingness to pay, linear incremental premium value, one provider and the maintained absence of capacity/quality choice, so it is a useful \emph{benchmark falsification} rather than a likely population constant.

\paragraph{Implication 3: local and cross-region returns to additional intermediation need not align.}
The project-level complementarity mechanism has a directly falsifiable component: at a given research effort $e$ and within the maintained interior success-probability domain, equation~\eqref{eq:success} implies
\begin{equation}
 \frac{\partial^2 p(e,\theta_A,\theta_B)}
 {\partial\theta_A\,\partial\theta_B}=\lambda\geq0.
 \label{eq:empirical-skill-complementarity}
\end{equation}
When $\lambda>0$, an increase in one partner's ex-ante technical capability raises the success-probability increment associated with the other partner's capability. This is a \emph{test of the maintained success technology}, not a novel theoretical theorem. It would require ex-ante skill measures for \emph{both} partners, recorded unsuccessful as well as successful collaborations, effort and project-characteristic controls, and a credible strategy for endogenous pairing and project selection. A positive observed interaction coefficient alone is not proof of causal complementarity.


Hold the neighbouring region $j$ at full founder participation, and compare region $i$'s private-only service with public/private coexistence under their respective equilibrium fee continuations. Lemma~3 implies that the neighbouring jurisdiction's additional surplus---which the local government financing region $i$ does not internalise---is
\begin{equation}
 \Delta W_j=
 C\left(S_j^0\ell_i+\frac{\lambda K_j}{4}\ell_i^2\right)
 =F_i^{\mathrm{national}}-F_i^{\mathrm{local}},
 \label{eq:empirical-spillover}
\end{equation}
where $K_j=M_j+\omega_j\Pi$. The sign of this \emph{foreign} benefit is not fixed: the consumer and founder returns from a new venture may be offset by business stealing from incumbent producers. The empirically discriminating question is therefore whether increases in cross-region matches or partner quality are accompanied by gains or losses \emph{where the new product competes}, not merely where its participating founders are registered.

A feasible data architecture would link eligible firms and their ex-ante skills, pair formation across jurisdictions, project success and product launch, product-market prices and quantities, incumbent performance, founder profit shares, programme budgets and tax-financed real costs. One must also estimate consumer surplus and the counterfactual outcomes of both markets \emph{without the venture}; aggregate collaboration statistics or participant turnover alone do not identify $M_j$ or either regional welfare threshold. Matched-pair projects and their sequential contracting milestones are observable in some intermediary datasets \citep{EldeboHjelm2025}, but their existence does not supply the necessary product-market counterfactual. Cross-border demand shocks, endogenous partner choice, market-entry selection and regional trade flows complicate identification.

\paragraph{A disciplined policy use.}
These implications motivate a diagnostic sequence, not an unestimated cost--benefit claim. First determine \emph{which} firms are reached and which technical skills they bring. Second distinguish basic-service access, private upgrade spending and supplier survival. Third identify where the commercialised products create consumer benefits, founder income and incumbent losses. Only after observing or credibly estimating these components could a government attempt the local-versus-aggregate welfare calculation. The narrow baseline ownership switch reported in Section~\ref{sec:numerical} can discipline a calibrated scenario analysis, but is \emph{not} a sign prediction for an average treatment effect of public hub openings. The model does not compare physical facilities with remote services, vouchers or procurement contracts, and none of these implications establishes that a new public building should be constructed.
